\documentclass[11pt]{article}
\usepackage{fontspec}
\usepackage[french]{babel}
\usepackage{amsmath}
\usepackage{amssymb}
\usepackage{amsthm}
\usepackage{mathrsfs}
\usepackage{tikz-cd}
\usetikzlibrary{babel}
\usepackage{enumitem}

\newcommand{\printedpage}[1]{\marginpar{\footnotesize p.~#1}}
\newcommand{\msp}[1]{\marginpar{\footnotesize [#1]}}
\newcommand{\D}{\mathbf{D}}
\newcommand{\cA}{\mathcal{A}}
\newcommand{\bDelta}{\boldsymbol{\Delta}}
\newcommand{\uHom}{\underline{\mathrm{Hom}}}

\begin{document}

\printedpage{}% no page number printed on this page
% head: (none)

\begin{center}
{\LARGE\textsf{LES DÉRIVATEURS}}

\bigskip

\textbf{ALEXANDRE GROTHENDIECK}

\bigskip\bigskip

\textbf{Chapitre IV}

\medskip

\textbf{Diagrammes substantiels}
\end{center}

\clearpage

\printedpage{}% no page number printed on this page
% head: (none)

\noindent Ce texte a été déchiffré et transcrit en \LaTeXe\ par M. Künzer. Il a été édité par M. Künzer, J. Malgoire, et G. Maltsiniotis. La transcription est aussi fidèle que possible au manuscrit. Pour les quelques corrections évidentes, ou rares commentaires des éditeurs, ainsi que pour la numérotation originale des pages du manuscrit, les caractères de machine à écrire \texttt{[typewriter]} entre crochets sont utilisés. Un point d'interrogation entre crochets signifie que l'on n'est pas sûr du mot qui précède.

\noindent Cette édition est provisoire. Les remarques, commentaires et corrections sont bienvenus. Envoyer un message à :

\noindent\textbf{maltsin@math.jussieu.fr}

\hfill G. Maltsiniotis\hspace*{2cm}

\clearpage

\printedpage{1}
% head: A. Grothendieck : Dérivateurs, Ch. IV, Édition des Universités de Montpellier II et Paris VII
% foot: Version 2 février 2002

\noindent\texttt{[Il n'y a pas de pages 1-72]}

\msp{page 73}

\subsection*{Construction de diagrammes substantiels pour un dérivateur.}

\noindent Soit $\D$ un prédérivateur, $I$ dans Diag, considérons

\noindent $\varphi_I : \mathcal{D}(I) \longrightarrow \uHom(I^o,\cA)$,\ \ où $\cA=\D(e)$ est la catégorie fondamentale des coefficients. % [sic] \mathcal{D}(I) -> \mathbf{D}(I), elsewhere the prederivator is the bold D

Les objets de $\uHom(I^o,\cA)$ s'interprètent comme des \textit{diagrammes de type $I^o$} dans $\cA$. Ce langage est surtout commode si $I$ est ordonnée. Notons, pour $F\in\mathrm{Ob}\,\D(I)$, par $(F_i)_{i\in\mathrm{Ob}\,I}$ le $I^o$-diagramme associé dans $\cA$. Par l'axiome Dér 2, le foncteur
\[
\varphi : F \longmapsto (F_i)_{i\in\mathrm{Ob}\,I}
\]
est conservatif. Mais il n'a aucune envie d'être fidèle, encore moins pleinement fidèle. Quand on connaît le diagramme $(F_i)_i$, on est loin de connaître $F$, pas même à isomorphisme non unique près, je présume. Il n'est nullement clair que si $F$, $F'$ dans $\D(I)$ sont tels que les diagrammes associés soient isomorphes, ils soient eux-mêmes isomorphes.

Pour insister sur ce point, les objets de $\D(I)$ s'appelleront les `\textit{diagrammes substantiels} de type $I^o$ dans $\cA$' (ou $\D$-substantiel, s'il y a risque de confusion). Pour étayer l'intuition et simplifier les explications, un objet $F$ de $\D(I)$ ($\D$-coefficients sur $I$) sera représenté graphiquement par le diagramme `faible' ou `squelettique' correspondant -- on le regarde
\msp{page 74}
comme une sorte de (pâle) `squelette' du diagramme substantiel dont il provient. P. ex. si $I=\bDelta_0 : 0\longrightarrow 1$ \texttt{[plutôt $\bDelta_1$]}, on note un diagramme substantiel de type $(\bDelta_0)^o$ \texttt{[plutôt $(\bDelta_1)^o$]} par le même graphisme, que le diagramme ordinaire :
\[
F_0 \longleftarrow F_1.
\]
Si
\[
I = \begin{tikzcd}
a \arrow[r] & c \\
0 \arrow[u] \arrow[r] & b, \arrow[u]
\end{tikzcd}
\]
les diagrammes substantiels de type $I$ se notent
\[
\begin{tikzcd}
F_a \arrow[d] & F_c \arrow[l] \arrow[d] \\
F_0 & F_b. \arrow[l]
\end{tikzcd}
\]

\clearpage

\printedpage{2}
% head: A. Grothendieck : Dérivateurs, Ch. IV, Édition des Universités de Montpellier II et Paris VII
% foot: Version 2 février 2002

Il faudra bien sûr préciser chaque fois si ce graphisme désigne un diagramme squelettique, ou substantiel. Souvent on part d'un squelette, et on essaye de `l'étoffer' avec un peu de chair, pour en faire un diagramme substantiel -- si possible, de façon canonique. C'est là un point de philosophie dérivatique essentiel pour la pratique. On procédera souvent de proche en proche, en `substantifiant' des sous-diagrammes correspondants à une suite strictement croissante $I_0\subset I_1\subset\cdots\subset I_n=I$ de $I$, au mieux de la situation particulière.

\msp{page 75}

En voici un exemple typique, en supposant $\D$ un dérivateur \textit{ponctué} (les $\D(I)$ sont ponctuées, et les $f^*$ transforment objets nuls en objets nuls). Il s'agit de la définition du foncteur `cosuspension' ou `espace de lacets' dans $\cA$,
\[
\Omega : \cA \longrightarrow \cA.
\]
Il est entendu que pour $M$ dans $\cA$, $\Omega(M)$ doit s'insérer dans un carré commutatif
\[
\begin{tikzcd}
e \arrow[d] & \Omega M \arrow[l] \arrow[d] \\
M & e, \arrow[l]
\end{tikzcd}
\tag*{(C)}
\]
où $e$ désigne les éléments unités. Ce carré cependant devrait être substantiel, plus précisément, être un carré substantiel `$\D$-\textit{cartésien}'. Étant entendu qu'un objet de $\D(I)$ (carré substantiel dans $\cA$) est dit $\D$-cartésien, s'il est dans l'image essentielle de $\D(I_1)$ par $i_{1*}$, où
\[
i_1 : I_1 \hookrightarrow I,
\]
$I_1$ étant la sous-catégorie ouverte de $I$ définie par
\[
I_1 = \begin{tikzcd}
a \\
0 \arrow[u] \arrow[r] & b.
\end{tikzcd}
\]

\msp{page 76}

La restriction du carré substantiel cherché (C) à la sous-catégorie $I_1$ est donc un diagramme substantiel de squelette
\[
\begin{tikzcd}
e \arrow[d] \\
M & e, \arrow[l]
\end{tikzcd}
\tag*{(Sq)}
\]
\printedpage{3}
% head: A. Grothendieck : Dérivateurs, Ch. IV, Édition des Universités de Montpellier II et Paris VII
% foot: Version 2 février 2002
dont il se déduit (à isomorphisme canonique près) en appliquant $i_*$ à ce dernier, de façon à le compléter en $\begin{tikzcd}
e \arrow[d] & \Omega M \arrow[l, dotted] \arrow[d, dotted] \\
M & e \arrow[l]
\end{tikzcd}$. Il reste donc à décrire une substantiation canonique du squelette (Sq), en termes de $M$. Cela nous conduit à regarder la sous-catégorie ouverte
\[
I_0 = \{0\} :\quad
\begin{tikzcd}
{} \arrow[r, dotted] & {} \\
0 \arrow[u, dotted] \arrow[r, dotted] & {} \arrow[u, dotted]
\end{tikzcd}
\]
de $I_1$ (donc aussi de $I$) réduite à l'objet initial, et l'inclusion
\[
I_0 \overset{i_0}{\hookrightarrow} I_1,
\]
et de substantifier le squelette (Sq) par $(i_0)_!(M)$.

On peut dire a posteriori que, si on connaissait le carré substantiel, il n'y a, à isomorphisme unique près, qu'une seule façon de l'obtenir comme un $i_{1*}i_{0!}(?)$, avec ? dans $\D(I_0)\simeq\cA$, car $i_{0!}$ et $i_{1*}$ étant pleinement fidèles, il en est de même de leur composé.

\msp{page 77}

En fait, on voit que le foncteur
\[
i_{1*}i_{0!} :\ \cA\simeq\D(I_0) \longrightarrow \D(I)
\]
induit une équivalence de $\cA$ avec la sous-catégorie pleine de $\D(I)$ formée des carrés $\D$-cartésiens dont les fibres en $a$ et $b$ sont nulles. (Le foncteur quasi-inverse évident est $i^*$, où $i : U_0 \hookrightarrow I$ \texttt{[plutôt $I_0 \hookrightarrow I$]}.) Cela signifie donc que dans ce cas, la donnée du squelette de nature très particulière ($F_a\simeq F_b\simeq 0_{\cA}$), plus la contrainte supplémentaire $\D$-cartésienne, sur le carré cherché, détermine celui-ci à isomorphisme unique près en termes de son squelette : il y a une substantiation non seulement `canonique' de celui-ci, mais même unique à isomorphisme unique près, à cause de la contrainte $\D$-cartésienne.

Je voudrais commencer par regarder le cas plus simple encore où $I=\bDelta_1=(0\longrightarrow 1)$, donc les squelettes envisagés sont simplement les flèches
\[
F_1 \longrightarrow F_0 \qquad \text{dans } \cA.
\]
La philosophie nous oblige, pour les flèches les plus fondamentales dans la théorie des dérivateurs $\D$ (flèches dans $\cA$, ou par extension dans une $\D(I)\simeq\cA_{\D_I}$, catégorie fondamentale du dérivateur induit $\D_I$) de se demander s'il n'y a pas

\end{document}
